% Oversells «Roman, nasr» shabloni (XeLaTeX yoki LuaLaTeX).
% Koʻrinish sahifa.sty da: muharrirdagi «Sahifa» tugmasi bilan uslublar va
% ustalarni oʻzgartiring. Matnni bob-bob alohida fayllarga ajratish mumkin:
% \include{boblar/bob1}.
\documentclass[11pt,twoside,openright]{book}
\usepackage[a5paper]{geometry}
\usepackage{fontspec}
\setmainfont{Noto Serif}
\usepackage{babel}
\babelprovide[import=uz-Latn, main]{uzbek}
\def\uzbekhyphenmins{22}
\usepackage{dizayn}
\usepackage{sahifa} % Oversells sahifa dizayni (sahifa.sty)

\begin{document}

\frontmatter
\sahifausta{muqaddima}
\begin{titlepage}
  \centering
  \vspace*{30mm}
  {\Large Muallif ismi\par}
  \vspace{20mm}
  {\Huge\bfseries Ona yurt\par}
  \vspace{6mm}
  {\large Qissa\par}
  \vfill
  {Toshkent — 2026\par}
\end{titlepage}

\chapter*{Soʻzboshi}
Bu shablon roman va qissalar uchun: boblar toq betda boshlanadi, bob oldidagi
boʻsh bet toza qoladi, juft betning yuqorisida kitob nomi, toq betning
yuqorisida bob nomi turadi.

\mainmatter
\sahifausta{asosiy}

\chapter{Qishloq tongi}
\begin{uslub}{boshi}
Tong otar-otmas qishloq uygʻondi. Tandirlardan koʻtarilgan tutun
terakzorlar ustida yupqa parda boʻlib yoyildi, ariq boʻyida yoʻlga
chiqqan sigirlarning qoʻngʻirogʻi eshitildi.
\end{uslub}

Otabek darvoza oldida toʻxtab, tepalikka qaradi. U yerda bolaligida
oʻynagan qari tut hamon turar, shoxlari yerga egilgan edi. Har yili bahor
kelganda bu tut birinchi boʻlib barg yozardi, qishloq ahli esa
\matnuslubi{urgu}{«tut uygʻondi — bahor keldi»} derdi.

\begin{uslub}{iqtibos}
Oʻz yurtingni sevmoq — oʻz onangni sevmoq kabidir.
\end{uslub}

Kun yoyilgach, dalaga chiqqanlar soni koʻpaydi. Kimdir ketmon koʻtargan,
kimdir arava haydar, bolalar esa ariq boʻyida qogʻoz qayiqchalar oqizardi.

\begin{uslub}{bezak}
* \quad * \quad *
\end{uslub}

\chapter{Yoʻl}
\begin{uslub}{boshi}
Shaharga olib boradigan yoʻl tor va changli edi. Otabek yelkasiga
tugunini olib, yoʻlga tushdi.
\end{uslub}

Yoʻl boʻyidagi har bir tosh, har bir quduq unga tanish edi. Bir paytlar
otasi bilan shu yoʻldan bozorga qatnagan, bozordan qaytishda doim
choyxona oldida toʻxtab, non va qovun yeyishardi.

\begin{uslub}{bezak}
* \quad * \quad *
\end{uslub}

\end{document}
